\chapter{Foundations: Conservation, Elements, and Models}
\label{ch:foundations}
\section{Why Circuit Theory Matters in Amateur Radio}
Amateur-radio equipment is full of circuits that select, transform, amplify,
measure, and deliver RF energy. A filter rejects an adjacent signal. A matching
network helps a transmitter see a reasonable load. A transmission line carries a
wave whose voltage and current may not be in phase with the source. Even when
the exam asks for a short formula, the underlying idea is usually conservation:
charge cannot disappear, energy must be accounted for, and stored energy creates
dynamics.

\section{KCL as an Accumulation Balance}
Kirchhoff's current law is a charge balance. At a node,
\[
\text{accumulation}=\text{inflow}-\text{outflow}.
\]
For an ordinary wire node with negligible stored charge, accumulation is zero,
so the algebraic sum of currents is zero:
\[
\sum_k i_k=0.
\]

\begin{controlsbox}
This is the same structure as a material balance in process systems. Capacitors
make the accumulation term visible because \(i_C=C\,\dd v_C/\dd t\).
\end{controlsbox}

\section{KVL as Energy Consistency}
Kirchhoff's voltage law says that the signed voltage rises and drops around a
closed loop sum to zero:
\[
\sum_k v_k=0.
\]
Voltage is energy per unit charge. KVL is therefore an energy-consistency rule:
after a charge goes around a closed loop and returns to the same point, its net
change in potential energy is zero.

\section{Constitutive Equations}
KCL and KVL tell us how elements connect. The element laws tell us what each
element does:
\[
v_R=Ri_R,\qquad i_C=C\frac{\dd v_C}{\dd t},\qquad
v_L=L\frac{\dd i_L}{\dd t}.
\]
Resistors dissipate energy. Capacitors store electric-field energy. Inductors
store magnetic-field energy:
\[
E_C=\frac{1}{2}Cv_C^2,\qquad
E_L=\frac{1}{2}Li_L^2.
\]

\section{State Variables}
The natural state variables are the quantities attached to stored energy:
capacitor voltages and inductor currents. A circuit with one independent energy
storage element is usually first order. A circuit with two independent storage
elements is usually second order.

\begin{physicalbox}
The physical circuit, the differential equation, the state model, and the phasor
model are not four different circuits. They are four coordinate systems for the
same object.
\end{physicalbox}

\section{Key Equations}
\[
\sum i_k=0,\qquad \sum v_k=0,\qquad
p(t)=v(t)i(t).
\]

\section{Key Definitions}
\begin{description}
\item[KCL] Conservation of charge at a node.
\item[KVL] Energy consistency around a loop.
\item[State] A variable needed, with the input, to predict future behavior.
\item[Passive sign convention] Positive absorbed power when current enters the
positive-voltage terminal.
\end{description}
